% ============================================================
% RSC Advances Introduction — ML/AI Emphasis with Application Context
% Flow: Agricultural/human-centric lighting need → TADF opportunity →
%       Computational bottleneck → Existing ML limitations →
%       Our contributions (accelerated discovery framework)
% ============================================================

\section{Introduction}\label{sec:introduction}

\subsection{Predicting the singlet--triplet gap of TADF emitters}

Thermally activated delayed fluorescence (TADF) emitters---organic donor--acceptor
molecules that harvest both singlet and triplet excitons for near-unity internal
quantum efficiency~\cite{Uoyama2012}---are the heavy-metal-free route to tunable,
efficient light sources for displays and for human-centric and horticultural
lighting~\cite{Yokoyama_2018, Bugbee2016}. An efficient emitter has to meet several
requirements at once: a small singlet--triplet gap
($\Delta E_\text{ST} < \qty{0.2}{\electronvolt}$), fast reverse intersystem crossing
($k_\text{RISC} > \qty{1e6}{\per\second}$), and emission matched to the target band.
This paper addresses only the first. The rate and device-efficiency requirements need
spin--orbit-coupling and device-level calculations beyond the cheap structural data
used here, and we make no quantitative claim about them (\cref{subsec:limitations}).

\subsection{The computational bottleneck and the role of machine learning}

Designing such emitters means searching a large space of donor--acceptor
combinations for a small $\Delta E_\text{ST}$. The direct route---time-dependent
DFT on excited states---scales steeply with system size~\cite{Curtarolo2013},
so exhaustive screening of more than a few hundred candidates is impractical.
Machine learning offers a cheaper surrogate~\cite{Butler2018, Pyzer-Knapp2022},
and several groups have predicted TADF gaps from molecular
descriptors~\cite{Gomez-Bombarelli2016, Barneschi2024}. Those low reported errors,
however, are obtained against \emph{computed} (TD-DFT or higher) reference gaps;
predicting the \emph{experimental} gap is markedly harder, to the point that one
study abandoned direct regression for classification~\cite{Nigam2024}
(\cref{subsec:comparison}).

Two problems recur in this literature, and we confront both. First, accuracy is
often quoted against a single random train--test split, which rewards
interpolation within familiar scaffolds and overstates how a model behaves on new
chemistry. Second, and more subtle, models are sometimes trained to reproduce a
\emph{computed} gap from a feature set that already contains the constituent
excitation energies; because $\Delta E_\text{ST}\equiv E_{S_1}-E_{T_1}$, this
reduces to subtraction and the reported accuracy reflects arithmetic rather than
learning.

Our earlier work set the stage: Article~1~\cite{Tchapet2024_xTB} benchmarked in JCIM
GFN2-xTB/sTDA against \num{747} experimental molecules
(MAE = \qty{0.162}{\electronvolt}), xTB geometry relaxations, and scaffold
cross-validation, and Article~2~\cite{Tchapet2024_DataDriven}
extracted structure--property design rules from that set. Here we ask a sharper question: what sets the ceiling on predicting the
\emph{measured} $\Delta E_{\text{ST}}$ from molecular structure, and what would it
take to raise it?

\subsection{Our contributions}

We do not claim a new high-accuracy predictor. The primary finding is that 2D
structure alone predicts the measured singlet--triplet gap as well as semi-empirical
excited-state features, because the experimental label, not the model, is the binding
constraint. Two further results are specific to this task and, to our knowledge, are
not reported elsewhere: the rank order that 2D structure encodes and the rank order
that only an excited-state calculation encodes come apart under a change of source
laboratory, and the feature--target leakage that inflates computed-gap models is
visible to some learners and hidden from others. Four points follow.

\paragraph{Structure alone predicts as well as the semi-empirical excited state.}
On \num{231} molecules across \num{212} Bemis--Murcko scaffolds, random forests built
from the 2D structure alone---Morgan fingerprints (MAE = \qty{0.091}{\electronvolt}) and
RDKit descriptors (\qty{0.100}{\electronvolt})---match descriptors derived from a
semi-empirical excited-state calculation (natural transition orbital, NTO, spatial
features, \qty{0.096}{\electronvolt}). The result holds across scaffold, cluster and random
cross-validation and against a label-permutation null, so it is not a splitting
artefact. On accuracy the excited-state step buys nothing: the two representations
agree to within \qty{0.017}{\electronvolt} on paired folds. On ranking under
laboratory transfer, however, the two representations separate. Under a DOI-grouped
split---holding out whole source publications rather than individual
molecules---the NTO model retains Spearman $\rho = \num{0.35}$ while Morgan drops to
\num{0.16} ($\Delta\rho = \num{0.19}$, \qty{95}{\percent} CI
\numrange{0.04}{0.34}, excluding zero only narrowly and on a single split, with
resampling over molecules rather than the \num{173} source papers). The excited-state
features do not improve accuracy on the gap; they preserve rank order that transfers
across source laboratories---a dissociation we have not seen reported for this
property. SHAP analysis ties the NTO
signal to hole--electron overlap and charge transfer, giving an interpretable read on
what structure encodes. Unlike a physics-guided model trained on NTO descriptors
against a \emph{computed} reference gap~\cite{Das2025}, we test the same class of
descriptor against the \emph{measured} gap and find the physical-guidance step buys
no accuracy there either.

\paragraph{The experimental noise floor, quantified.}
Why does accuracy plateau near \qty{0.1}{\electronvolt} regardless of descriptor set?
Independent experimental reports for the same compound scatter by
\qty{0.072}{\electronvolt} RMS (the worst molecule reaching a standard deviation of
\qty{0.56}{\electronvolt}), and the median gap the model is trained on is itself known
only to \qty{0.049}{\electronvolt}---about half the model error. That
$\Delta E_\text{ST}$ is reported inconsistently has been shown for individual
emitters~\cite{Shuaib2026}; what we add is the size of that inconsistency over
\num{231} molecules and \num{173} source papers and its direct consequence for
prediction---to our knowledge the first inter-laboratory variability budget for the
TADF singlet--triplet gap. The label bounds what any predictor can reach, consistent
with the broader finding that many machine-learning benchmarks in drug-, materials-
and molecular-discovery are limited by dataset noise rather than model
capacity~\cite{Crusius2025}---applying their $R^2_\text{max} = 1-\sigma^2_\text{noise}/\sigma^2_y$
framework to our label distribution gives $R^2_\text{max} \approx 0.71$--\num{0.91}
depending on which noise floor applies (\cref{subsec:ceilings})---and with the aleatoric/epistemic distinction under
which label noise is an irreducible floor no model can cross~\cite{Heid2023}.
Computed references face a separate limit: B3LYP and
CAM-B3LYP gaps differ by up to \qty{0.58}{\electronvolt}, so a gap computed at one
level of theory does not transfer to another; that bound constrains the meaning of a
computed reference, not the error of a model fitted to measured labels. The
\qty{0.1}{\electronvolt} plateau is not a model-capacity failure: neither a
higher-capacity tabular learner (gradient boosting, MLP, SVR) nor a frozen pretrained
molecular-transformer embedding beats the random forest (the latter reaches only MAE
\qtyrange{0.109}{0.113}{\electronvolt}), and the result places itself honestly at
$R^2 \approx 0.25$, a real but limited signal.

\paragraph{Feature--target leakage, and why it survives review.}
A feature that is a proxy for the target---here, an exact one---is a long-recognised
cause of over-optimistic ML results~\cite{Kaufman2012,Kapoor2023,Artrith2021}, and the
apparent severity of such a leak is known to depend on the model class. What we add is
the controlled measurement for this specific case.
Regressing $\Delta E_\text{ST}$ on features containing the $S_1$ and $T_1$ energies is
circular ($\Delta E_\text{ST} \equiv E_{S_1} - E_{T_1}$): under the scaffold
cross-validation used throughout, a linear model given those energies recovers the gap
exactly ($R^2 = 1.000$, MAE $= 0$) against $R^2 = 0.36$ without them. An axis-aligned
random forest does not recover that subtraction---its splits approximate a difference
of two continuous inputs rather than express it, a familiar limitation of tree
ensembles---so it reaches only $R^2 = 0.59$ on the identical leaked features, a score
that does not look anomalous, which is how a tree-based study can carry this defect
through cross-validation unremarked. The lesson
transfers to any computational-chemistry ML pipeline that regresses on computed energy
components.

\paragraph{Modest triage utility, and what would be needed to push past the ceiling.}
As a pre-screen, the model enriches low-gap candidates
(\numrange{1.2}{1.4}$\times$ over random selection) on held-out labelled data, though
a bootstrap confidence interval shows this only holds at cutoffs of \num{50}
molecules or more---at the ten-to-twenty-candidate scale a laboratory would actually
shortlist, the enrichment is not distinguishable from chance. On
the filtered virtual library the conformal prediction intervals are too wide to support
molecule-level selection; no ranked shortlist is deposited. $\Delta E_\text{ST}$ is
necessary but not sufficient for TADF efficiency, and we make no device-level
performance claims. Progress past the $\approx\qty{0.1}{\electronvolt}$ ceiling requires
not a better model but a better dataset: a curated, single-source, solvent-controlled
experimental benchmark of sufficient size to push the label precision below the present
\qty{0.049}{\electronvolt} floor. This study provides the quantitative justification
for that investment.

This study builds on our earlier work---\textbf{Article~1}~\cite{Tchapet2024_xTB}
(xTB/sTDA validation) and \textbf{Article~2}~\cite{Tchapet2024_DataDriven} (design
rules)---by measuring, rigorously, what such cheap data can and cannot support.
All data, trained models, and code are openly available (Zenodo DOI:
10.5281/zenodo.17436069), enabling reproduction and extension of the benchmark.
